\section{Discussion}

Our investigation into automated mathematical discovery systems yields several important insights regarding the fundamental nature of novel mathematical concept generation. The experimental results on the Collatz sequences provide an illuminating case study that demonstrates both the potential and limitations of current approaches.

The observed pattern irregularity in the Collatz sequences, particularly the lack of correlation between input size and sequence properties, suggests a broader principle about mathematical discovery: the most interesting mathematical structures often exhibit behavior that is simultaneously structured yet unpredictable. This aligns with the theoretical framework developed in Section 2, where we showed that for any recursive function $f: \mathbb{N} \to \mathbb{N}$, the associated $_f$ sequence space exhibits similar properties of structured unpredictability.

A key finding from our analysis is that the space of potential mathematical concepts appears to be hierarchically structured. Let $\mathcal{M}_n$ denote the space of mathematical concepts of complexity $\leq n$ (under our proposed measure). Our results suggest that:

\[
\frac{|\mathcal{M}_{n+1} \setminus \mathcal{M}_n|}{|\mathcal{M}_n|} \geq c \cdot n
\]

for some constant $c > 0$, indicating that the space of possible concepts grows super-exponentially with complexity. This has important implications for automated discovery systems, as it suggests that simple enumeration-based approaches will necessarily miss most interesting structures.

The experimental results align well with recent work by \cite{tsoukalas2025learning} on learning mathematical interestingness, though our approach differs in focusing on structural properties rather than heuristic measures. Our findings complement those of \cite{georgiev2025mathematical}, who demonstrated similar hierarchical organization in their large-scale mathematical exploration system.

A limitation of our current approach is that it may not capture certain types of mathematical concepts that resist formalization in our framework. This is particularly evident in the case of geometric intuitions, which often precede formal algebraic descriptions. Future work could explore extending our model to incorporate geometric and topological primitives, perhaps building on the automated reasoning tools described in \cite{zarzuelo2025experience}.

The implications for automated mathematical discovery are significant: while our results show that systematic exploration of mathematical concepts is possible, they also suggest that purely automated approaches may need to be complemented by human mathematical intuition for optimal results. This hybrid approach, combining machine learning with mathematical expertise, appears to be the most promising direction for future research in this area.